\chapter{Guía de evaluación}
\label{ch:guia}

Los capítulos anteriores especifican qué evalúa el marco LLM-Sec: un modelo de madurez de seis niveles, una taxonomía de 62 controles y tres perfiles organizacionales que orientan la interpretación del resultado. Este capítulo especifica cómo se conduce la evaluación. Define el flujo de siete pasos que va desde la selección del perfil hasta el plan de remediación, los criterios que el evaluador aplica para dar por cumplido un control según su fuente de evidencia, la matriz de responsabilidades que la organización debe instanciar para que la evaluación sea auditable, y el ciclo iterativo mediante el cual se cierran las brechas identificadas.

\section{Flujo de evaluación}
\label{sec:flujo-evaluacion}

El proceso de evaluación del marco LLM-Sec sigue siete pasos secuenciales:

\begin{enumerate}[leftmargin=2em]
  \item \textbf{Selección del perfil organizacional:} identificar el perfil (P1, P2 o P3) que corresponde a la organización según su tipo y obligaciones regulatorias.
  \item \textbf{Evaluación de controles:} verificar cada control de las subdimensiones D1 a D5 utilizando la fuente de evidencia correspondiente (configuración, experimental o declarativa).
  \item \textbf{Cálculo de métricas:} ejecutar la auditoría ofensiva y calcular TFA, LAS, PCC y NCP conforme al capítulo~\ref{ch:metricas}. Estas métricas constituyen la evidencia experimental que sustenta la verificación de los controles del paso anterior.
  \item \textbf{Asignación de nivel por subdimensión:} el nivel de cada subdimensión es el mayor nivel cuyos controles están completamente implementados, respetando la progresión acumulativa.
  \item \textbf{Cálculo del nivel global:} aplicar la regla de piso mínimo: $L = \min(d_1, d_2, d_3, d_4, d_5)$.
  \item \textbf{Análisis de brechas:} identificar la distancia entre el nivel verificado y el nivel objetivo del perfil seleccionado.
  \item \textbf{Plan de remediación:} priorizar la implementación de controles faltantes según el nivel objetivo, enfocándose primero en la subdimensión que actúa como piso.
\end{enumerate}

% TODO: Insertar diagrama del flujo de evaluación (Figura 3.13 de la tesis, adaptada).

\section{Criterios de verificación por fuente de evidencia}
\label{sec:criterios-verificacion}

\subsection{Configuración}

El control se verifica inspeccionando la declaración explícita en los archivos de configuración del sistema. Ejemplos: presencia de umbrales en el archivo YAML de la capa defensiva, instrucciones de seguridad en el prompt del sistema, particionamiento de la base vectorial.

\subsection{Experimental}

El control se verifica mediante la ejecución de pruebas de auditoría ofensiva contra el sistema bajo evaluación. El resultado se expresa como una métrica cuantificable (tasa de filtrado, latencia, cobertura). Las pruebas ejecutadas deben cubrir los vectores OWASP declarados como fuente normativa del control: un control anclado a LLM02:2025 no puede verificarse con pruebas que solo ejerciten inyección de instrucciones.

\subsection{Declarativo}

El control se verifica mediante revisión documental o entrevista. El evaluador solicita evidencia de la existencia de la política, procedimiento o asignación de roles descrita por el control.


\section{Matriz RACI}
\label{sec:raci}

El marco define roles genéricos que cada organización debe mapear a personas o cargos concretos. La matriz RACI establece las responsabilidades para las actividades principales del ciclo de evaluación.

La gestión efectiva de la seguridad de un sistema RAG requiere que cada actividad tenga un responsable claramente definido. Siguiendo el modelo de COBIT 4.1, el marco LLM-Sec provee una matriz RACI como artefacto central de gobernanza, que define los roles genéricos involucrados en la operación del sistema junto con su nivel de participación en cada actividad.

Los roles definidos son genéricos por diseño: no corresponden a personas específicas sino a funciones organizacionales. Cada organización debe instanciar la matriz asignando cada rol a los cargos o personas concretas de su estructura, conforme al control C50. Esta instanciación es condición necesaria para que los controles de D5 sean auditables.

La Tabla~\ref{tab:raci} presenta las actividades principales del marco y la asignación de roles según los criterios R (Responsable), A (\textit{Accountable}), C (Consultado) e I (Informado).

\begin{table}[H]
    \centering
    \caption{Matriz RACI del marco LLM-Sec.}
    \label{tab:raci}
    \renewcommand{\arraystretch}{1.4}
    \begin{tabular}{
        >{\raggedright\arraybackslash}p{5cm}
        >{\centering\arraybackslash}p{1.2cm}
        >{\centering\arraybackslash}p{1.2cm}
        >{\centering\arraybackslash}p{1.2cm}
        >{\centering\arraybackslash}p{1.2cm}
        >{\centering\arraybackslash}p{1.2cm}
    }
    \toprule
    \textbf{Actividad} &
    \rotatebox{45}{\textbf{Resp. Seguridad}} &
    \rotatebox{45}{\textbf{Operador Técnico}} &
    \rotatebox{45}{\textbf{Prop. Corpus}} &
    \rotatebox{45}{\textbf{Dirección}} &
    \rotatebox{45}{\textbf{Cumplimiento}} \\
    \midrule
    Designar responsable de seguridad del sistema RAG & A & I & I & R & C \\
    Declarar el sistema RAG como activo de información & R & C & C & A & I \\
    Mantener inventario de activos actualizado & R & C & C & I & A \\
    Gestionar accesos al sistema & R & A & C & I & C \\
    Aplicar lista de chequeo LLM-Sec & R & A & C & I & C \\
    Revisar y actualizar el marco periódicamente & R & C & C & A & C \\
    Notificar incidentes de seguridad & R & C & I & A & C \\
    Ejecutar plan de pruebas ofensivas & A & R & I & I & C \\
    Capacitar a los roles en seguridad LLM & R & C & C & A & C \\
    Reportar métricas de gobernanza a la dirección & R & C & I & A & C \\
    \bottomrule
    \end{tabular}
    \begin{minipage}{\linewidth}
        \vspace{4pt}
        \footnotesize \textbf{R} = Responsable (ejecuta) \enspace\textbullet\enspace
        \textbf{A} = \textit{Accountable} (rinde cuentas) \enspace\textbullet\enspace
        \textbf{C} = Consultado \enspace\textbullet\enspace
        \textbf{I} = Informado
    \end{minipage}
\end{table}

\section{Ciclo de mejora}
\label{sec:ciclo-mejora}

Ante brechas identificadas entre el nivel verificado y el nivel objetivo, el evaluador debe:

\begin{enumerate}[leftmargin=2em]
  \item Identificar la subdimensión que actúa como piso mínimo.
  \item Priorizar los controles faltantes de esa subdimensión en el nivel inmediatamente superior.
  \item Implementar los controles priorizados.
  \item Re-ejecutar la auditoría ofensiva y recalcular las métricas del capítulo~\ref{ch:metricas}: la TFA verifica que la implementación produjo mejora defensiva efectiva, la LAS caracteriza el costo operacional introducido, el PCC confirma la cobertura de controles alcanzada y el NCP verifica que la mejora es consistente entre vectores de ataque.
  \item Re-ejecutar la evaluación para verificar el tránsito de nivel.
\end{enumerate}

Este ciclo es iterativo y se repite hasta alcanzar el nivel objetivo del perfil organizacional o hasta que los riesgos residuales sean formalmente aceptados y documentados por la dirección.